%% kapa: exam=1 round=36 year=2025 period=2 form=A subject=accounting subject_ko=회계학 seq=40 number=80
%% source: Q-Net 감정평가사 기출문제 (cst00309) artlSeq=5231525 file=raw/1008/5231525/2025년 제36회 감정평가사 1차시험 2교시 문제지 원본.pdf
%% conversion: kordoc 4.24.51 -> tools/gen1.py (auto); review=proofread-sonnet; flags=none
\begin{question}{80}
전략적 관리회계기법에 관한 내용으로 옳은 것을 모두 고른 것은?
\begin{framedtext}
ㄱ. 제약이론(theory of constraints)은 기업의 목표를 달성하는 과정에서 병목 공정을 파악하여 집중적으로 관리하고 개선해서 기업의 성과를 높이는 방법이다. \par
ㄴ. 품질원가(cost of quality)에서 예방원가와 평가원가를 포함하는 통제원가는 불량품의 발생률과 역의 관계를 갖는다. \par
ㄷ. 활동기준경영(activity based management)에서 활동분석을 통하여 파악된 비부가가치활동은 검사, 이동, 대기 및 저장 등의 활동이 있다. \par
ㄹ. 적시생산시스템(just-in-time)은 적정한 안전재고를 보유하기 위한 기법으로, 초변동원가계산을 적용하여 제품원가를 계산한다. \par
ㅁ. 제품수명주기원가계산(life cycle costing)은 제품의 제조단계에 초점을 맞춰 제품설계에 따라 제품을 생산하면서 지속적인 개선을 통하여 원가절감 방안을 모색한다.
\end{framedtext}
\begin{choices}
  \item ㄱ, ㄴ, ㄷ
  \item ㄱ, ㄷ, ㅁ
  \item ㄱ, ㄹ, ㅁ
  \item ㄴ, ㄷ, ㄹ
  \item ㄴ, ㄹ, ㅁ
\end{choices}
\end{question}
\answer{1}
